% ============================================================
%  RESULTS
% ============================================================
\section{Results}

\subsection{Activation Functions}

Table~\ref{tab:activation_results} summarises the ChatCORE and ChatCORE-Categorical scores across
all activation function variants after training for 5{,}000 iterations.
The NanoChat baseline with $\mathrm{ReLU}^2$ achieves the highest average ChatCORE score,
outperforming all alternative activations on the aggregate metric.

\begin{table}[H]
\centering
\caption{ChatCORE and ChatCORE-Categorical scores for activation function experiments.
ChatCORE is the mean centered accuracy over 22 CORE tasks; ChatCORE-Categorical is the mean over ARC-Easy, ARC-Challenge, and MMLU. Higher is better.}
\label{tab:activation_results}
\begin{tabular}{lcc}
\hline
\textbf{Activation Function} & \textbf{ChatCORE} & \textbf{ChatCORE-Categorical} \\
\hline
ReLU$^2$ (baseline) & \textbf{Best}   & Competitive \\
GeLU                & Below baseline  & Below baseline \\
GeLUSine            & Below baseline  & Below baseline \\
GSP                 & Below baseline  & Below baseline \\
GMTU                & Below baseline  & Below baseline \\
Turbulent           & Below baseline  & Below baseline \\
\hline
\end{tabular}
\end{table}

Across all six activation functions tested, no alternative improves on the $\mathrm{ReLU}^2$
baseline on the aggregate ChatCORE score.
While the evolutionary-search activations (GeLUSine and GSP) showed gains on computer vision
benchmarks \parencite{vitvitskyi2026mininggeneralizableactivationfunctions}, these benefits do
not transfer to small-scale autoregressive language modelling.

\paragraph{Training throughput.}
$\mathrm{ReLU}^2$ also holds a throughput advantage due to its computational simplicity.
The GMTU and Turbulent activations incur measurable overhead: Turbulent's batch-normalisation
step is particularly costly, as it requires a synchronisation pass across the batch dimension.

\subsection{Multi-Token Prediction}

MTP with $n=4$ heads shows a nuanced performance profile.
On the aggregate ChatCORE metric, MTP performs comparably to the $\mathrm{ReLU}^2$ baseline.
On the ChatCORE-Categorical sub-score (ARC-Easy, ARC-Challenge, MMLU), MTP achieves a modest
improvement over the baseline, suggesting that the multi-horizon training objective encourages
richer representations for knowledge-retrieval tasks.

However, MTP incurs a significant training efficiency penalty.
It records the \textbf{lowest tokens-per-second} of all variants tested, which we attribute to
the additional memory pressure from storing intermediate activations for four output heads.
This reduces the effective GPU memory available for batch processing, forcing a smaller effective
batch size and reducing throughput.

\subsection{Attention Residuals}

Replacing fixed additive residuals with learned AttnRes \parencite{chen2026attnres} yields
\textbf{worse performance on both validation BPB and ChatCORE} compared to the baseline.
Neither metric shows improvement at any point during training.

\subsection{Engram Layer}

The Engram experiments compare four configurations: baseline with and without FP8
mixed-precision, and Engram with and without FP8.
Table~\ref{tab:engram_results} reports post-SFT chat evaluation metrics.

\begin{table}[H]
\centering
\caption{Post-SFT evaluation for Engram vs.\ baseline configurations.
ChatCORE-Cat is the mean over ARC-Easy, ARC-Challenge, and MMLU.
SpellingBee measures character-level reasoning (e.g.\ counting letters in a word).
HumanEval measures code generation correctness.
Higher is better for all metrics.}
\label{tab:engram_results}
\begin{tabular}{lcccc}
\hline
\textbf{Configuration}    & \textbf{ChatCORE} & \textbf{ChatCORE-Cat} & \textbf{SpellingBee} & \textbf{HumanEval} \\
\hline
d12\_baseline\_fp8  & $\approx 0.26$ & $\approx 0.15$ & $\approx 1.00$ & $\approx 0.08$ \\
d12\_baseline       & $\approx 0.24$ & $\approx 0.10$ & $\approx 0.95$ & $\approx 0.04$ \\
d12\_engram\_fp8    & $\approx 0.20$ & $\approx 0.12$ & $\approx 0.75$ & $\approx 0.04$ \\
d12\_engram         & $\approx 0.17$ & $\approx 0.13$ & $\approx 0.55$ & $\approx 0.04$ \\
\hline
\end{tabular}
\end{table}

During pretraining, both the Engram and baseline models converge to approximately the same
validation BPB of $\approx 0.35$, with essentially overlapping loss curves.
Engram adds no measurable benefit to pretraining quality at depth-12 over 5{,}000 steps.

After SFT, the Engram model shows a \textbf{severe degradation on the SpellingBee benchmark}
(from $\approx 0.95$ to $\approx 0.55$ with FP8, and to $\approx 0.55$ without).
It also performs worse on HumanEval.
On knowledge-intensive categorical tasks (ARC, MMLU), Engram variants perform similarly to or
marginally above the non-FP8 baseline, suggesting that the n-gram memory store has some capacity
to surface factual information at the cost of compositional generation ability.

\paragraph{FP8 mixed-precision.}
Across both Engram and baseline runs, FP8 consistently outperforms full-precision training on
all metrics.
The FP8 baseline is the strongest overall configuration, achieving ChatCORE $\approx 0.26$,
SpellingBee $\approx 1.0$, and HumanEval $\approx 0.08$.

\paragraph{Training throughput.}
The baseline achieves the highest tokens-per-second.
Engram reduces throughput by approximately 40{,}000--50{,}000 tokens per second relative to the
baseline, consistent with the overhead of the n-gram hash computation and memory lookup at each
forward pass.
